\section{Implications}\label{sec:implications}

\textbf{Universities and skilling providers.}
\begin{itemize}
  \item Make \emph{statistics and probability} and \emph{one statistical language plus SQL} the
        non-negotiable core; SAS Viya for Learners removes the licence cost of the SAS track.
  \item Add a graded \emph{data storytelling} module (dashboards, visual narrative, presenting to
        stakeholders). It is the most under-served skill: rewarded by employers, rarely taught or
        advertised.
  \item Teach the big-data stack (Hadoop, Spark, Hive, Scala) as one bundle, once trainer capacity
        allows.
  \item Plan trainer capacity before budget: sharing trainers across institutions unlocks the next
        courses more cheaply than extra money.
\end{itemize}

\textbf{Students.} Skills linked to pay and growth are concrete and learnable: statistics, a
programming toolkit, and the ability to explain results. Postings that centre on MIS reporting sit
in lower salary bands; building analysis and modelling skills moves candidates up.

\textbf{Employers and HR.} Storytelling and quantitative depth predict who earns high hikes, yet
job postings rarely ask for storytelling. Writing it into job descriptions and assessments would
align hiring with what employers already reward. For senior, customer-facing roles, structured
assessment of conscientiousness and openness, alongside mentoring, is supported by both this data
and prior research \citep{barrick1991}.

\textbf{Policy and CSR programmes.} The same model can be run for a district or a group of
institutions to decide where shared trainers and seats create the most value, which suits CSR-funded
skilling where budgets are fixed and capacity is scarce.

\subsection*{Limitations}
\begin{itemize}
  \item All results are \textbf{associations} from cross-sectional sample data, not causal effects.
  \item JDS and SDS are \textbf{small} (139 and 161); intervals are wide and SDS separates unusually
        cleanly. Validation on another cohort is needed before policy use.
  \item Salary is available only in \textbf{bands}; descriptions are truncated at source; 42\% of tag
        uses are data-derived skills without an ESCO/O*NET code.
  \item RQ4 course costs, seats, hours and coverage are \textbf{documented assumptions}; with higher
        per-seat costs (our first draft, before the source check) budget rather than trainers became
        binding. The conclusion depends on cost levels, so they are stated explicitly and are easy to
        change.
\end{itemize}
